\documentclass[11pt,a4paper]{article}
\usepackage[T1]{fontenc}\usepackage[utf8]{inputenc}
\usepackage[margin=2.0cm]{geometry}
\usepackage[expansion=false]{microtype}
\usepackage{enumitem,amsmath,amssymb,booktabs,titlesec,amsthm}
\usepackage[hidelinks]{hyperref}
\setlist[itemize]{leftmargin=1.1em,itemsep=1pt,topsep=2pt}
\setlength{\parskip}{2pt}\emergencystretch=2em
\titleformat{\section}{\large\bfseries}{\thesection.}{0.4em}{}
\titlespacing*{\section}{0pt}{8pt plus 1pt}{3pt}
\newtheorem{prop}{Proposition}
\begin{document}
\begin{center}
{\large\bfseries The circle map of the sleep switch in closed form}\\[2pt]
{\small Appendix to the report on the falsification programme. 30 October 2026.}
\end{center}
\vspace{-2pt}

\noindent\textbf{Summary.} The sleep--wake layer of the model reduces \emph{exactly} to a relay with
hysteresis, and its return map admits a closed form up to two transcendental equations. This yields:
the folds of the fast subsystem in explicit parametric form ($1.4503$ and $2.4635$\,mV against
$1.450/2.463$ in the model); the free-running period in closed form; the edge of the $1{:}1$ tongue
as the solution of a system of three equations, with no scan. It also settles the question of
invariant tori: there are \emph{none} in the window, and the proof is an inequality on the Lyapunov
exponent rather than a parameter scan. A by-product is that the Arnold-tongue framework does not
apply to this system at all: the map is piecewise smooth and does not cover the circle.

\section{Reduction to a relay with hysteresis}
The fast pair of nucleus voltages ($\tau\approx10$\,s) is eliminated adiabatically against the
somnogen time scales (18.2 and 4.2\,h). Its steady state admits a \emph{parametric} solution: taking
$V_m$ as the parameter,
\begin{equation}\label{eq:par}
Q(V_v)=\frac{A_m-V_m}{|\nu_{mv}|},\qquad V_v=Q^{-1}\!\Bigl(\frac{A_m-V_m}{|\nu_{mv}|}\Bigr),\qquad
D_v(V_m)=V_v-\nu_{vm}Q(V_m),
\end{equation}
with $Q^{-1}(q)=\theta+\sigma_Q\ln\bigl(q/(Q_{\max}-q)\bigr)$. The folds are the extrema of
$D_v(V_m)$. Evaluating~\eqref{eq:par}: $D_-=1.4503$\,mV (sleep branch, $Q(V_m)=0.75$) and
$D_+=2.4635$\,mV (wake branch, $Q(V_m)=3.45$), against the $1.450$ and $2.463$ quoted in the model
as a check against the source publication. The agreement to four significant figures is obtained
\emph{without integration}.

In the limit $\tau\to0$ the system is a relay: the somnogens obey linear equations with closed-form
solutions, and switching occurs at two \emph{moving thresholds}
\begin{equation}\label{eq:thr}
g_\pm(t)=\frac{D_\pm-B_0+\kappa-\nu_{vc}\cos\bigl(2\pi(t-\alpha)/24\bigr)}{A_G},
\qquad \kappa=k\,D .
\end{equation}

\begin{prop}[Return map]\label{prop:map}
Let $t_n$ be a sleep onset. Then $t_{n+1}=f(t_n)$ is given implicitly by
\begin{equation}\label{eq:map}
g_+(t_n)\,e^{-(t_w-t_n)/\tau_s}=g_-(t_w),\qquad
1-\bigl(1-g_-(t_w)\bigr)e^{-(t_{n+1}-t_w)/\tau_w}=g_+(t_{n+1}),
\end{equation}
where $t_w$ is the wake onset. Both equations are of the form ``exponential $=$ sinusoid''. The map
is smooth wherever the crossing is transversal and loses smoothness where it is tangential; the
latter condition is the critical locus.
\end{prop}

\section{Free period and the delayed-bifurcation correction}
For $\nu_{vc}=0$ the thresholds are constant and~\eqref{eq:map} is solved explicitly:
\begin{equation}\label{eq:per}
T_s=\tau_s\ln\frac{g_+}{g_-},\qquad
T_w=\tau_w\ln\frac{1-g_-}{1-g_+},\qquad T_0=T_s+T_w .
\end{equation}
Substitution gives $T_s=1.337$\,h, $T_w=6.683$\,h, $T_0=8.020$\,h. Direct integration at zero
amplitude gives $8.878$\,h, a discrepancy of 10\,\%. The cause is identified and is not an error of
derivation: switching does \emph{not} occur at the fold. Sleep onset is observed at $D_v=2.5026$
against the fold at $2.4635$, and wake onset at $1.3794$ against $1.4503$: the system overshoots the
fold, because the relaxation time diverges at a saddle-node and a slow passage completes after the
bifurcation point. The overshoot was measured over two decades in the fast time scale and vanishes
with exponent $0.64$ and $0.59$, against the value $2/3$ for slow passage through a saddle-node.
Hence~\eqref{eq:per} is the \emph{exact} $\tau\to0$ limit and the model stands 10\,\% away from it;
when analytic thresholds are used this offset must be stated rather than concealed.

\section{The 1:1 tongue edge without a scan}
The $1{:}1$ orbit is a fixed point of $f$ shifted by one day. The tongue edge is its saddle-node:
\begin{equation}\label{eq:edge}
f(t^*)=t^*+24,\qquad f'(t^*)=1,
\end{equation}
which together with~\eqref{eq:map} gives three equations in $(t^*,t_w,\kappa)$ and determines
$\kappa_{\mathrm{edge}}$ to machine precision. Numerical bisection on a locking criterion gave
$4.7651$\,mV; an independent scan of the rotation number (rule~F20) gave $4.77$\,mV; the difference
is $0.005$\,mV.

\section{No invariant tori: the exponent is bounded away from zero}
\textbf{The map is not defined everywhere.} The somnogens rise towards unity from below, so sleep
onset is possible only where $g_+(t)<1$. The fraction of the day satisfying this is
\begin{center}\small
\begin{tabular}{lccccc}
\toprule
$\kappa$, mV & 0 & 2 & 4.80 & 5.10 & 5.18\\
fraction of the day admitting sleep onset & 0.789 & 0.531 & 0.166 & 0.077 & 0.016\\
\bottomrule
\end{tabular}
\end{center}
At $\kappa\approx5.18$ this fraction vanishes, which is how sleep is lost. At $\kappa=4.80$ the
domain is the arc $[10.014,\,13.986]$\,h of length $3.97$, and the image is the arc
$[10.846,\,11.839]$ of length $0.993$, contained in the domain.

It follows immediately that the Arnold-tongue framework \textbf{does not apply}. A degree-one circle
map has an image covering the whole circle; here the image occupies 4\,\% of it. The map is
piecewise smooth with discontinuities: at $\kappa=4.88$ the jump is $22.47$\,h against a median step
of $10^{-4}$\,h, that is, almost exactly one day. This is a skipped night. \emph{The locking
staircase is generated by branch switching, not by mode locking of a smooth circle map.}

\textbf{The correct quantity is the exponent along the orbit.} The maximum of $|f'|$ over a grid is
not suitable: at tangential crossings the derivative has a square-root singularity, and a grid
samples points the attractor never visits. Iterating~\eqref{eq:map} and discarding the transient:
\begin{center}\small
\begin{tabular}{lcccccc}
\toprule
$\kappa$, mV & 4.78 & 4.82 & 4.88 & 4.94 & 5.00 & 5.10\\
$\lambda$ along the orbit & $-3.17$ & $-3.01$ & $-3.24$ & $-3.34$ & $-3.54$ & $-2.53$\\
rotation number $\rho$, days & 1.0000 & 1.5000 & 2.0000 & 2.0000 & 2.0000 & 2.0000\\
\bottomrule
\end{tabular}
\end{center}
Over the whole window $\lambda$ lies between $-1.62$ and $-3.54$. Two conclusions follow, both as
inequalities rather than observations. \emph{Chaos is excluded}: it requires $\lambda>0$. \emph{A
smooth invariant torus is excluded}: a quasiperiodic attractor requires $\lambda=0$ exactly, whereas
here the distance to zero is at least $1.6$. The attractor at every parameter examined is a strongly
attracting periodic orbit. The rotation numbers $1$, $3/2$, $2$ reproduce the staircase of rule~F20,
which is a third independent confirmation of the reduction, after the folds and after the tongue
edge.

\textbf{What a grid could not do.} In its first edition, part~(zh) of rule~F30 asserted that a
locked fraction below unity on a grid implies that tori occur on a parameter set of positive
measure. The inference is invalid. The complement of the locked set is \emph{dense}, so a finite
grid finds unlocked nodes whether the complement has positive measure or measure zero, and misses
them in both cases with a probability depending only on the resolution: the criterion has no power
in either direction. The honest content of the outcome ``240 of 240 nodes locked'' is a lower bound
on the measure of the locked set, $0.988$ at 95\,\% confidence, and nothing more.

Reaching the correct quantity required discarding two incorrect ones: first the locked fraction on a
grid, then the maximum of $|f'|$ over the domain. The general lesson: on a Poincar\'e section, a
quantity averaged \emph{along the orbit} and a quantity taken over the \emph{whole domain} answer
different questions, and the latter need not concern the attractor at all.

\textbf{Measure in phase space is a different question and does not arise here.} All of the above
concerns \emph{parameter} space. The measure of the attractor itself --- a thick attractor of
positive Lebesgue measure, whose dimension equals that of the system --- belongs to a different and
substantially harder circle of problems (Ilyashenko constructed examples for step skew products).
The present system sits at the opposite pole, and this is measured rather than assumed. The fast
time scales (10\,s) against the slow ones (hours) give a strongly negative divergence, and the two
fast dimensions collapse immediately: control~F30(a) established one-dimensionality of the return
map to $4.5\times10^{-13}$, eight trajectories from different initial somnogen values falling onto a
single curve to machine precision. The one-dimensional map then contracts onto a periodic orbit with
$\lambda\le-1.6$, so the attractor on the section is a finite set of points and, in the extended
phase space, a closed curve. This is structural: approaching thickness would require the divergence
to cease being negative somewhere, that is, the fast time scales to become comparable to the slow
ones --- but the entire reduction to a relay, and with it the folds and the circle map, rests on
their separation.

\section{Consequences for computation}
Constructing the map from~\eqref{eq:map} requires two scalar roots per cycle instead of 1600
integration steps per day, and the tongue edge from~\eqref{eq:edge} requires solving three equations
instead of a bisection. The F30 run took 1686\,s; the analytic route yields the same numbers in
seconds, up to the 10\,\% threshold offset stated above. This applies in particular to part~(e),
which the run returned as untestable: the exponent along the orbit is obtained in twelve iterations
of the map and answers its question directly, whereas a grid of 240 nodes could not answer it in
principle, since it searched for aperiodic nodes where the attractor is periodic at every parameter.

\smallskip\noindent\textbf{Open points.} (i)~The overshoot exponent converged to $0.59$--$0.64$
against the theoretical $2/3$; the asymptotic regime was not reached, and either a further decade in
the fast time scale is needed or the correction contains a logarithm. (ii)~The exponent along the
orbit was computed at twelve values of $\kappa$; formally this does not exclude exceptional
parameters carrying a Denjoy-type Cantor minimal attractor. The argument against is structural
rather than statistical: such an attractor also requires an exponent at zero, and the exponent
nowhere rises above $-1.6$. (iii)~The separation of branches of the map was performed by thresholding
the jump size and is not reliable; the number of branches is deliberately not quoted. (iv)~All of the
above concerns the deterministic switch with frozen damage; with live damage and load noise the
question of attractors does not arise in this form.
\end{document}
